% a promotional effort to show that the workflow does work and
% produce good results

% keep source lines to <= 79 characters
%234567891_234567892_234567893_234567894_234567895_234567896_234567897_23456789

% bio:
% Timothy Prescott (https://campus.und.edu/directory/timothy.prescott) is an associate professor in the department of Mathematics & Statistics at the University of North Dakota, where he maintains the department's Calculus textbook and coordinates the data for the emporium math lab.

\DocumentMetadata{
  lang=en,
  tagging=on,
  pdfstandard=ua-2,
%  check-tagging-status,
  tagging-setup={
%    math/alt/use,
    math/setup=mathml-SE
  }
}

\title[pdftitle={Converting a PDF textbook to be accessible}
]{Converting a \PDF\ textbook to be accessible}
\author{Timothy Prescott}
\date{}

\documentclass{article}
\usepackage{unicode-math}
\usepackage{microtype}
\usepackage{hyperref}

\makeatletter
\newcommand{\SMC}{%
  \ifx\@currsize\normalsize\small\else
   \ifx\@currsize\small\footnotesize\else
    \ifx\@currsize\footnotesize\scriptsize\else
     \ifx\@currsize\large\normalsize\else
      \ifx\@currsize\Large\large\else
       \ifx\@currsize\LARGE\Large\else
        \ifx\@currsize\scriptsize\tiny\else
         \ifx\@currsize\tiny\tiny\else
          \ifx\@currsize\huge\LARGE\else
           \ifx\@currsize\Huge\huge\else
            \small\SMC@unknown@warning
 \fi\fi\fi\fi\fi\fi\fi\fi\fi\fi
}
\makeatother

\newcommand*{\acro}[1]{{\SMC #1}\@}
\newcommand{\PDF}{\acro{PDF}}
\newcommand{\veraPDF}{vera\PDF}
\newcommand{\VeraPDF}{vera\PDF}  % it appears they are always lower case
\newcommand{\PDFix}{\PDF ix}

\setcounter{secnumdepth}{-1}

%\address{University of North Dakota}
%\netaddress{timothy dot prescott (at) UND dot edu}
%\personalURL{https://und.edu/directory/timothy.prescott}
%\ORCID{0009-0009-0555-8074}

\begin{document}
\maketitle

\begin{abstract}
Modern \LaTeX\ provides a viable roadmap to creating accessible \PDF s.
%Convincing others of that fact may be a different matter.
%Creating a large accessible \PDF\ is straightforward using modern \LaTeX.
%Convincing others you have done so may be a different matter.
\end{abstract}

\section{Introduction \& setup}

For the past decade, the University of North Dakota has used a
\href{https://arts-sciences.und.edu/academics/math/calc-1-texts.html}{custom 1150 page textbook}
for its three semester calculus sequence.
This has enabled students
to download the \PDF\ for free
in addition to having the option to purchase the book
for a fraction of the cost of other textbooks
(currently \$25--\$30 per semester, when I last checked).

My university has been emphasizing the national requirement
that all documents must satisfy the \acro{WCAG 2.1 AA} guideline
as of April 2027.
Beginning in early 2025, I began exploring how to ensure that
the textbook meets this guideline.

\section{The basics}

Because we create the textbook using \LaTeX,
we were able to leverage much of the work of the \LaTeX\ project.
By placing a single command at the beginning of the code's preamble,
we indicated that Lua\LaTeX\ should produce
a \PDF\ version 2.0 satisfying the \acro{UA-2} standard.

Because tagging in \LaTeX\ is still under development,
I did encounter several bugs in the resulting compilation.
These were generally quickly fixed or circumvented by users
in online forums (often members of the \LaTeX\ project).

\section{Alt text and tables}

A key part of any accessibility workflow is providing alt text for images.
With Lua\LaTeX\ automatically handling the bulk of \PDF\ tagging,
this was the greatest part of the remaining work.
%providing alternative text for the thousand or so images throughout the book.
\LaTeX\ can specify alt text
to the image, and also provides the capability to note that an image
is an artifact or actual text.

It is useful to note that creating alt text is
not necessarily a \TeX nical aspect of the tagging process,
but possibly the most time intensive.
For alt text, I used a short Python script that added
\texttt{alt=\{ALT-TEXT-TO-BE-DETERMINED\}} to all of the images.
A student worker with less \TeX\ experience was then able to go through
the book, examine the resulting \PDF, and help create the alt text.

There are two contrasting problems we encountered with alt text.
Alt text guidelines generally prefer a description that is too short
to adequately describe a complicated mathematical object.
I typically considered but not necessarily followed such a restriction.

A bigger problem arises in early chapters in the calculus sequence
when we ask students to verbally describe a graph.
It is not clear how we can provide alt text to describe a graph
without giving an answer that students could repeat.
I chose to describe such graphs using various synonyms
(``the graph is a line from $(-2,0.5)$ to a hollow dot at $(-1,1)$, \dots''),
but acknowledge this is not an entirely satisfactory approach.

Many \PDF\ checkers do not recognize \PDF\ version 2.0 and
instead (silently) check the \PDF\ against version 1.7.
A critical distinction for us is that \PDF\ version 1.7 requires that
formulas include alt text.
We can accomplish this with \LaTeX,
but this usually causes \PDF\ screen readers to read the alt text and
omit the more informative Math\acro{ML} tags
(despite a recent recommendation to the contrary),
%\url{https://pdfa.org/resource/best-practice-guide-math-in-pdf/},
so I discourage its use
unless a document is being judged on its adherence to the 1.7 standard.

The last remaining significant accessibility task is that
tables need to identify the headings that describe the data in the table.
But we should emphasize that if the sole purpose of the table
is to arrange elements in a grid,
then it's not actually a table and probably shouldn't be tagged as such.
Again, \LaTeX\ can identify the headers,
although the current syntax to do so is still under development.

\section{Checking the results}

Having resolved the reported errors,
I needed to verify that the result was indeed accessible.
The obvious first step is to ensure
a screen reader properly handles the document.
Unfortunately, when it comes to \PDF\ version 2.0,
the only currently viable \PDF\ readers are Adobe, Foxit, and Firefox,
while the only currently viable screen readers to go with them are
\acro{NVDA} and \acro{JAWS} on Windows and Orca on \acro{GNU}/Linux.

Another approach to verifying accessibility is to verify that
the \PDF\ is indeed version 2.0,
and satisfies the \acro{UA-2} \PDF\ standard.
Unfortunately (again),
there are not many \PDF\ checkers that recognize version 2.0,
and many will apply version 1.7 standards to a version 2.0 document,
resulting in many claimed ``issues'' with a valid document.
In the positive direction,
%\url{https://pdfa.org/accessible-math-in-pdf-finally/}
there are several (but currently only
five, and at least three of the five are different user interfaces
over the same validation profiles)
checkers that support \PDF\ version 2.0.
% listed at \url{https://pdfa.org/accessible-math-in-pdf-finally}.
I have experience with
\href{https://verapdf.org/}{\veraPDF} and
\href{https://pdfix.net/}{\PDFix}.

\hypersetup{hidelinks}

The \LaTeX\ team uses \veraPDF\ to validate their documents.
\VeraPDF\ also has a command line version of the program,
allowing for faster, or automated, checking of generated \PDF s.
After a minute of checking, \veraPDF\ reported that the textbook
satisfies the \acro{UA-2} and \acro{WTPDF} standards,
having performed over eleven million checks
of one thousand seven hundred rules.

\PDFix\ uses a similar checker as \veraPDF, but it also provides a
graphical view of the \PDF\ in question, which better locates errors.
\PDFix\ also allows navigating the tree structure of the \PDF,
making it possible to locate errors given by \veraPDF.
Additionally, \PDFix\ offers additional profiles for validation.

My institution uses Ally to check \PDF s posted to our \acro{LMS}.
Ally is one of the aforementioned \PDF\ checkers that looks for version 1.7
and requires formulas to have alt text.
For this reason, I try to avoid posting the textbook to our \acro{LMS}.

%For shorter documents, I can also examine the tagging structure
%of a \PDF\ by using the command line tool \texttt{show-pdf-tags}
%or ng\acro{PDF} at \url{https://ngpdf.com/}.

\section{Other textbooks and lessons learned}

Since converting the calculus textbook to be accessible,
I have also converted both a linear algebra textbook
% (see \url{https://github.com/teepeemm/LinearAlgebra})
and a differential equations textbook
% (see \url{https://github.com/teepeemm/trenchDEs})
to be accessible.

%Converting these textbooks creates
%a more readable, more maintainable, and more accessible document.
%I would argue that the one-time cost of doing so
%is well worth the few hours required for the result.

I have encountered some misconceptions during this process.
Some people think that \PDF s, especially those produced by \LaTeX,
are inherently inaccessible.
usually because they are not aware that
this reality has completely changed over the past few years.
A separate misconception is that inaccessibilities in a \LaTeX\ \PDF\ must
be laboriously remediated elsewhere (usually in Adobe Pro),
not realizing that \LaTeX\ has the ability
to create a fully accessible \PDF.
\bigskip

%As a final remark, I have collected the advice from this article
%into a GitHub repository at\\
%%\cite{accessibleURL}.
%\url{https://teepeemm.github.io/accessibility}.

A longer version of this article is available at
\url{https://tug.org/TUGboat/tb47-2/tb146prescott-book-accessibility.pdf}.

\end{document}
